\begin{problem}
    Let $f:\mathbf{R}^n\to\mathbf{R}$ be Lebesgue measurable, and let $g:\mathbf{R}^n\to\mathbf{R}$
    be a function which agrees with $f$ outside of a set of measure zero, thus there exists a set
    $A\subseteq\mathbf{R}^n$ of measure zero such that $f(x)=g(x)$ for all $x\in\mathbf{R}^n\setminus A$.
    Show that $g$ is also Lebesgue measurable.
    (\emph{Hint:} use Exercise~7.4.10.)
\end{problem}

\begin{problem}
    Suppose that \(f:\mathbf{R}\to [0,\infty)\) is Lebesgue measurable and
    \(g:[0,\infty)\to [0,\infty)\) is monotone. Prove that
    \(g\circ f\) is Lebesgue measurable.
\end{problem}

\begin{proof}
    From the claim in the proof of Lemma 7.5.5, we have
    \[
    (g \circ f)^{-1}((a, \infty)) = f^{-1}(g^{-1}((a, \infty))).
    \]
    Let
    \[
    I_a := g^{-1}((a, \infty)) = \{y \in [0, \infty) : g(y) > a\}.
    \]
    Since $g$ is monotone defined on $[0, \infty)$, $I_a$ can only be one of the following forms for some $c \ge 0$:
    \[
    \emptyset, \quad [0, \infty), \quad [0, c), \quad [0, c], \quad (c, \infty), \quad [c, \infty).
    \]
    We evaluate the measurability of $f^{-1}(I_a)$:
    \begin{itemize}
        \item If $I_a = (c, \infty)$: $f^{-1}(I_a)$ is measurable by definition.
        \item If $I_a = [c, \infty)$: $f^{-1}(I_a)= \bigcap_{n=1}^{\infty} f^{-1}((c - 1/n, \infty))$. Since the countable intersection of measurable sets is measurable, $f^{-1}(I_a)$ is measurable.
        \item If $I_a = [0, c) \text{ or } [0, c]$: Since 
        \[
        f^{-1}([0, c))=\mathbb{R} \setminus f^{-1}([c,\infty)) \quad \text{and} \quad f^{-1}([0, c])=\mathbb{R} \setminus f^{-1}((c,\infty)),
        \]
        $f^{-1}(I_a)$ is measurable.
        \item If $I_a = \emptyset \text{ or } [0, \infty)$: $f^{-1}(I_a)$ are $\emptyset$ and $\mathbb{R}$ respectively, so both of them are measurable.
    \end{itemize}
    Since $f^{-1}(I_a)$ is measurable for all possible forms of $I_a$, $(g \circ f)^{-1}((a, \infty))$ is measurable for every $a \in \mathbb{R}$. Hence, by Lemma 7.5.8, $g \circ f$ is Lebesgue measurable.    
\end{proof}

\begin{problem}
    A function \(f:\mathbf{R}\to \mathbf{R}\) is \emph{upper semicontinuous} if
    \[
    \lim_{k\to\infty} x_k=x
    \quad\Longrightarrow\quad
    \limsup_{k\to\infty} f(x_k)\le f(x).
    \]
    
    
    \begin{enumerate}
        \item[(a)] Draw a graph of an upper semicontinuous function that is not continuous.
        
        \item[(b)] Show that upper semicontinuity is equivalent to the requirement that for every open ray \((-\infty,a)\), the preimage
        \[
        f^{-1}((-\infty,a))
        \]
        is an open set.
         \item[(c)] Show that an upper semicontinuous function is measurable.
        
    \end{enumerate}
\end{problem}

\begin{proof}[(a)]
    Suppose $f : \mathbb{R} \to \mathbb{R}$ is
    \[
    f(x)=
    \begin{cases}
    1,& x=0,\\
    0,& x\neq 0.
    \end{cases}
    \]
    \begin{center}
    \begin{tikzpicture}[scale=1.2]
        % axes
        \draw[->] (-3,0) -- (3,0) node[right] {$x$};
        \draw[->] (0,-0.5) -- (0,1.8) node[above] {$y$};
    
        % horizontal line y=0
        \draw[thick] (-3,0) -- (3,0);
    
        % open circle at (0,0)
        \draw[fill=white, thick] (0,0) circle (2pt);
    
        % filled dot at (0,1)
        \fill (0,1) circle (2pt);
    
        % labels
        \node[left] at (0,1) {$1$};
        \node[below left] at (0,0) {$0$};
    \end{tikzpicture}
    \end{center}

    $f$ is not continuous because
    \[
        0 =\lim_{x \to 0} f(x) \neq f(0) = 1.
    \]

    Because $f(x) \le 1 $ for all $x \in \mathbb{R}$, so if $x_k \to 0$, $\limsup_{k\to\infty} f(x_k)\le 1 = f(0)$ is true. 
    
    If $x_k \to a$ and $a \ne 0$, then by definition we can pick $N$ large enough such that $\vert x_n - a \vert \le \frac{a}{2}$ for all $n  \ge N$, and hence $f(x_n) = 0$ for all $n \ge N$. Then we can conclude that 
    \[
    \limsup_{k\to\infty} f(x_k)\le 0 = f(a).
    \]

    So we can conclude that $f$ is not continuous and it is semicontinuous.
\end{proof}

\begin{proof}[(b)]
    \vphantom{text}
    \begin{itemize}
        \item [\((\implies )\)] Let $U_a = \{x \in \mathbb{R} : f(x) \in (-\infty, a)\}$ and we want to show that $U_a$ is open for every $a$.

        For fixed $a$, suppose $U_a$ is not open, then there exist some point $y \in U_a$ such that $y$ is not the interior point of $U_a$. Then we can construct a sequence $(y_n)_{n=1}^{\infty}$ such that
        \[
            y_n \notin U_a \text{ and } \lvert y_n - y \rvert < \frac{1}{n}.
        \]
        It is obvious that $y_n \to y$, and since $f$ is semicontinous, 
        \[
            \limsup_{k\to\infty} f(y_k)\le f(y).
        \]
        However since $y \in U_a$, $f(y) < a$. And since $y_k \notin U_a$ for all $k$, 
        \[
            \limsup_{k\to\infty} f(y_k) \ge a.
        \]
        This lead to a contradiction, and hence $U_a$ is an open set.
        \item [\((\impliedby )\)] Suppose $x_k \to x$, we want to show that 
        \[
            \limsup_{k\to\infty} f(x_k) \le f(x).
        \]
        It suffices to show that for every $\varepsilon > 0$,
        \[
            \limsup_{k\to\infty} f(x_k) \le f(x) + \varepsilon.
        \]
        For fixed $\varepsilon > 0$, let $a = f(x) + \varepsilon$. Since $f(x) < a$, $x \in f^{-1}((-\infty, a))$. And since $f^{-1}((-\infty, a))$ is an open set, $x$ is the interior point of $f^{-1}((-\infty, a))$. There exist some $\delta > 0$ such that
        \[
            \lvert y -x \rvert < \delta \implies y \in f^{-1}((-\infty, a)).
        \]
        And since $x_n \to x$, we may pick $N$ large enough such that for all $n \ge N$, $\lvert x_n - x \rvert < \delta$. Then $x_n \in f^{-1}((-\infty, a))$ for all $n \ge N$.

        This show that 
        \[
            \limsup_{k\to\infty} f(x_k) < a = f(x) + \varepsilon,
        \]
        and hence we have proved.
    \end{itemize}
\end{proof}

\begin{proof}[(c)]
    It suffices to show $f^{-1}((-\infty, a))$ is measurable for all $a$. Because if we know this is true, then for any interval $(a,b)$, 
    \[
        \begin{aligned}
            (a, b) &= (-\infty, b) \cap (\mathbb{R} \setminus (-\infty, a]) \\
            &= (-\infty, b) \cap \bigcap_{n=1}^{\infty} (-\infty, a + \frac{1}{n}).
        \end{aligned}
    \]
    Therefore, 
    \[
        f^{-1}((a,b)) = f^{-1}((-\infty, a)) \cap \bigcap_{n=1}^{\infty} f^{-1}((-\infty, a + \frac{1}{n})),
    \]
    where $f^{-1}((a,b))$ is the intersection of countably many measurable set, and hence it is measurable.

    For any open set, it can be written as the union of countably many open intervals, by similar argument, we can conclude that the preimage of open set is measurable.

    By previous subproblem, we know that for semicontinuous function $f$, $f^{-1}((-\infty, a))$ is an open set, and hence it is measurable by Borel property. And we are done.
\end{proof}

\begin{problem}
    Let $\Omega$ be a measurable subset of $\mathbb{R}^n$, and let $f:\Omega\to[0,\infty]$ and $g:\Omega\to[0,\infty]$ be measurable functions.
    Without using Theorem 8.2.9 or Lemma 8.2.10, prove that
    $$
    \int_\Omega (f+g)\ge \int_\Omega f + \int_\Omega g.
    $$
\end{problem}

\begin{proof}
    Let $\phi, \psi$ be nonnegative simple function such that
    \[
        0 \le \phi \le f, 0 \le \psi \le g.
    \]
    Then 
    \[
        0 \le \phi + \psi \le f + g.
    \]
    Also $\phi + \psi$ is a simple function, therefore by definition,
    \[
        \int_\Omega (f+g) \ge \int_\Omega (\phi + \psi),
    \]
    and since the integration of simple function is additive,
    \[
        \int_\Omega (f+g) \ge \int_\Omega \phi + \int_\Omega \psi.
    \]
    And hence
    \[
        \int_\Omega (f+g) \ge \sup_{\phi \le f, \psi \le g} (\int_\Omega \phi + \int_\Omega \psi) = \sup_{\phi \le f} \int_\Omega \phi + \sup_{\psi \le g} \int_\Omega \psi.
    \]
    By definition, this implies that
    \[
        \int_\Omega (f+g) \ge \int_\Omega f + \int_\Omega g.
    \]
    And we are done.
\end{proof}

\begin{problem}
    Let \(\{s_n\}\) be an increasing sequence of nonnegative simple functions which converges pointwise on an interval \(I\) to a limit function \(f\). If \(I\) is unbounded and if \(f(x)\ge 1\) almost everywhere on \(I\), prove that the sequence
    \[
    \left\{\int_I s_n\right\}
    \] diverges.
\end{problem}

\begin{proof}
    Since each $s_n$ is measurable and the sequence converges pointwise to $f$, by Lemma 7.5.10, $f$ is also measurable. 
    
    Let $A=\{x \in I: f(x) \ge 1\}$. Since $f(x) \ge 1$ almost everywhere on $I$, we have 
    \[
    m(I \setminus A)=0.
    \]
    
    This implies $m(A)=\infty$ since $I$ is unbounded. For each integer $n \ge 1$, define
    \[
    E_n := \{x \in A: s_n(x) > \frac{1}{2}\}.
    \]
    
    Since $s_n$ is a simple function, it is measurable, which makes each $E_n$ a measurable set.
    
    Additionally, since $\{s_n\}$ is an increasing sequence, we have
    \[
    E_1 \subseteq E_2 \subseteq E_3 \subseteq \dots .
    \]
    
    Choose any $x \in A$, we have 
    \[
    f(x) \ge 1 > \frac{1}{2}.
    \]
    
    Since $\lim_{n \to \infty} s_n(x) = f(x)$, there exists N such that for every $n \ge N$,
    \[
    s_n(x) > \frac{1}{2}.
    \]

    Hence, every $x \in A$ belongs to $E_n$, and 
    \[
    A = \bigcup_{n=1}^{\infty} E_n.
    \]

    By the continuity of Lebesgue measure, we have 
    \[
    m(A) = \lim_{n \to \infty} m(E_n).
    \]
    
    Since $m(A)= \infty$, we get 
    \[
    \lim_{n \to \infty} m(E_n) = \infty.
    \]

    Let $\chi_{E_n}$ denote the characteristic function of $E_n$.
    
    Since $s_n(x)>\frac{1}{2}$ on $E_n$ and $s_n \ge 0$, we have
    \[
    s_n \ge s_n \chi_{E_n} \ge \frac{1}{2} \chi_{E_n}.
    \]
    
    By Prop 8.1.10(d), we get
    \[
    \int_I s_n \ge \int_I \frac{1}{2} \chi_{E_n}.
    \]
    
    And by Lemma 8.1.9 and Prop 8.1.10(c), we get
    \[
    \int_I s_n \ge \int_I \frac{1}{2} \chi_{E_n} = \frac{1}{2} \int_I \chi_{E_n} = \frac{1}{2} m(E_n).
    \]
    
    Taking $n \to \infty$ on both sides, we can conclude that $\lim_{n \to \infty} \int_I s_n = \infty$. Hence, $\{\int_I s_n\}$ diverges.    
\end{proof}

\begin{problem}
    Let \(\{s_n\}_{n=1}^\infty\) be a sequence of measurable simple functions on \(\mathbf{R}\) such that
    \[
    s_n(x)\to f(x)\qquad \text{for every }x\in \mathbf{R}.
    \]
    \begin{enumerate}
        \item[(a)] Prove that for every \(a\in \mathbf{R}\),
        \[
        f^{-1}\bigl((a,\infty)\bigr)
        =
        \bigcup_{n=1}^\infty \bigcap_{k=n}^\infty
        s_k^{-1}\!\left(\left(a+\frac1n,\infty\right)\right).
        \]
        
        \item[(b)] Deduce that \(f\) is measurable.
    \end{enumerate}
\end{problem}